\bookchapter{Chapter 27: Making Phones Easy for Parents — Complete Senior-Friendly Guide}

\chapterintro{Honestly, when writing this chapter, I thought of my own parents. Every time a “system update” pops up on my mom's phone, she's afraid to tap it, worried she'll break the phone. My dad's WeChat font is smaller than an ant, but he doesn't know how to enlarge it — he's been squinting at it for two years. It's not that they're “dumb” — it's that phones default to assuming everyone is a twenty-something. In this chapter, I use the plainest language to solve the most frustrating problems elderly users face with their phones, one by one.}

\booksection{First, Get the Mindset Right: It's Not “Seniors Can't Use Phones” — It's “Phones Weren't Designed for Them”}

Before making any settings changes, one concept needs to be clear:\textbf{Current smartphone interfaces were never really designed with elderly users in mind}. Pop-ups are too small to see, gestures require precise touches, icons all look the same — these aren't user problems, they're product problems. So when helping elderly relatives with phones, the first thing isn't to teach — it's to help them\textbf{set up the environment}. Once the environment is set up properly, most problems disappear.

Three iron rules for helping elderly relatives with phones: (1)\textbf{Don't rush}— you've used smartphones for ten years, they may have used them for three months. (2)\textbf{Use analogies}— “Wi-Fi is like the water pipes at home — turn on the faucet and water flows.” This kind of metaphor is infinitely easier to understand than technical jargon. (3)\textbf{Write it down}— they might forget verbal instructions. Write the key steps on a piece of paper and stick it to the back of their phone case — it works better than anything else.

\booksection{I. First Step: Enable “Easy Mode” — The Phone Becomes Clean and Simple}

Most phones now have a built-in “Easy Mode” (some call it “Care Mode,” “Senior Mode,” or “Elder Mode”). After enabling it, the entire phone transforms:

\begin{itemize}
  \item \textbf{Fonts and icons become very large}— no more need for reading glasses to see the screen.
  \item \textbf{Desktop icons greatly simplified}— only Phone, Messages, Camera, Photos, WeChat, and a few other commonly used apps remain; everything else is tucked away.
  \item \textbf{Operations become simpler}— many flashy features (leftmost page, smart assistant, floating windows, etc.) are disabled by default — the interface is clean and straightforward.
  \item \textbf{Sound becomes louder}— ringtone, notification, and call volumes are automatically set to maximum.
\end{itemize}

\booksubsection{How to enable Easy Mode by brand (just follow along):}

\begin{center}
\begin{tabularx}{\linewidth}{|X|X|X|}
\hline
\textbf{Brand} & \textbf{Mode Name} & \textbf{How to Enable} \tabularnewline
\hline
iPhone & Accessibility (no separate Easy Mode, but items can be enlarged individually) & Settings > Display \& Brightness > Text Size, drag to max. Then go to "Accessibility > Display \& Text Size" and enable "Bold Text" and "Button Shapes" \tabularnewline
4. Wireless Charging Protocols & Easy Mode & Settings > System \& Updates > Easy Mode > Enable \tabularnewline
Xiaomi/Redmi & Minimalist Mode & Settings > Additional Settings > Minimalist Mode > Enable \tabularnewline
OPPO/OnePlus & Easy Mode & Settings > Special Features > Easy Mode \tabularnewline
vivo/iQOO & Easy Desktop & Settings > Home Screen, Lock Screen \& Wallpaper > Easy Desktop \tabularnewline
Samsung & Easy Mode & Settings > Display > Easy Mode \tabularnewline
\hline
\end{tabularx}
\end{center}

If you use another brand, search for “simple” or “elder” in the search box at the top of Settings — you'll usually find it. If not, follow the “Manual Enlargement” approach below step by step — the result is the same.

\booksection{II. If You Don't Use Easy Mode, Manual Adjustment Works Too — “Big Text + Big Icons” in Five Steps}

Some elderly relatives don't want to use Easy Mode — they think it looks too “old-fashioned.” In that case, manually adjust these five items for the same effect:

\begin{enumerate}
  \item \textbf{Set Font to Maximum}: Settings > Display > Font Size, drag to the far right. Also set "Display Size" (sometimes called "Display Zoom") to maximum — this makes all interface elements larger together.
  \item \textbf{Make Icons Bigger}: Pinch on the home screen (like zooming out on a photo) to enter desktop settings > Desktop Layout (or Icon Size), and choose the layout with the fewest grid cells (like 3×5) — icons naturally become bigger.
  \item \textbf{Enable "High Contrast Text"}: Makes text edges sharper — find this toggle in "Accessibility."
  \item \textbf{Disable Unnecessary Notifications}: Settings > Notifications, turn off all app notifications except Phone, Messages, and WeChat. This way the phone won't buzz every few minutes and won't feel annoying.
  \item \textbf{Put Frequently Contacted People on the Home Screen}: Long-press empty space on home screen > Add Widget > Contacts, and place speed-dial shortcuts for spouse, children, doctor, and emergency contacts on the home screen. One tap to call — no need to scroll through contacts.
\end{enumerate}

\begin{tcbpractice}{Hands-On: Help Elderly Relatives "Renovate" Their Phone}
\begin{enumerate}
  \item Enable Easy Mode (follow the brand table above).
  \item Turn on battery percentage display (Settings > Battery > Battery Percentage), so they always know how much charge is left.
  \item Open the notification shade and remove or hide quick toggles like Bluetooth, NFC, and Airplane Mode that elderly relatives don't need.
  \item Place four speed-dial contacts on the home screen: spouse, children, family doctor (if any), community/property management.
\end{enumerate}

\end{tcbpractice}

\booksection{III. Anti-Fraud — Please Read This Section Together with Your Elderly Relatives}

Elderly people are prime targets for telecom fraud. It's not that they're less vigilant — scammers' scripts are carefully designed to exploit elderly people's psychological vulnerabilities. Here are the most common scam types and how to respond. I recommend reading them aloud to your elderly relatives, or printing them out and sticking them next to the phone.

\booksubsection{Six Common Scams (Hang Up Immediately If You Encounter These):}

\begin{itemize}
  \item \textbf{(1) "Your Grandson Is in Trouble, Send Money Now"}— The scammer says "your grandson/granddaughter got into trouble outside, was arrested by police, and needs bail." Remember:\textbf{Any call demanding a transfer to "release someone" is a scam}. Police and government agencies never demand money over the phone.
  \item \textbf{(2) "Your Health Insurance Card Has Been Suspended"}— A text message says your health insurance card is abnormal and asks you to click a link to "activate" it. Remember:\textbf{Health insurance and social security bureaus never send text messages asking you to click links}. In this situation, call 12333 (the official social security hotline) to verify.
  \item \textbf{(3) "High-Yield Investment, Guaranteed Profit"}— WeChat groups and text messages recommend various "retirement investments" and "government-supported projects," promising high returns. Remember:\textbf{Anything promising annualized returns over 5\% that's "guaranteed" is 99.9\% a scam}.
  \item \textbf{(4) "I'm Your Boss/Old Friend, I Changed My Number"}— Impersonating an acquaintance to borrow money. Remember:\textbf{Calling to verify works better than anything}. If they say "I can't take calls," they're almost certainly a scammer.
  \item \textbf{(5) "Free Eggs / Health Product Seminar"}— Free offline events lure elderly people into attending "health seminars" and then sell overpriced health products. Remember:\textbf{There's no free lunch — free events are always businesses in disguise}.
  \item \textbf{(6) "Your Phone Is Infected, Needs Remote Repair"}— The scammer asks you to download a remote control app (like Sunflower, TeamViewer), then remotely controls your phone to transfer money. Remember:\textbf{Absolutely never let anyone remotely control your phone!}Phone manufacturers, carriers, and banks will never call you asking you to install remote control software.
\end{itemize}

\booksubsection{Three defense lines that, once set up, can block 90\% of scams:}

\begin{enumerate}
  \item \textbf{First Line: Built-in Phone Blocking}— On Android, enable "Harassment Blocking" (in the Phone app or Phone Manager app). On iPhone, enable "Silence Unknown Callers" (Settings > Phone > Silence Unknown Callers).
  \item \textbf{Second Line: Carrier Blocking}— Send text "KTFSR" to 10086 to activate China Mobile's free high-frequency spam call blocking. China Telecom and China Unicom have similar services. Call your carrier's customer service to have them enable it for you.
  \item \textbf{Third Line: National Anti-Fraud Center App}— Download the "National Anti-Fraud Center" app from the app store, register, and enable incoming call alerts. This app is linked to police databases — scam numbers are identified as soon as they call.
\end{enumerate}

\begin{tcbwarn}{Anti-Fraud Mantra for Elderly Relatives (Print and Stick Next to the Phone)}
\begin{itemize}
  \item \textbf{Stranger calls asking for money transfer Hang up!}
  \item \textbf{Text message links asking you to click Delete!}
  \item \textbf{Promising high returns with guaranteed principal Block!}
  \item \textbf{Asking you to install remote control software Hang up immediately!}
  \item \textbf{Real police don't handle cases over the phone, and definitely don't ask for money!}
  \item \textbf{When unsure about anything, call your children first.}
\end{itemize}

\end{tcbwarn}

\booksection{IV. Remote Assistance — Fix Your Parents' Phone from Anywhere}

Sometimes elderly relatives encounter problems that you can't explain over the phone — they can't find the right button. That's when “remote assistance” comes in — you can see their phone screen on your own phone and help them operate it. Each brand has a different solution:

\begin{itemize}
  \item \textbf{4. Wireless Charging Protocols}: Built-in "Changlian" app has a "Screen Sharing" feature. Both you and your parents install Changlian, start a video call, and tap "Share Screen" — you can draw circles on their screen to point out "tap here."
  \item \textbf{Xiaomi/Redmi}: Xiaomi Video Call app has "Remote Assistance," supporting screen sharing + simultaneous voice call.
  \item \textbf{[Install Apps from External Sources]}: Built-in "Family Space" or "Remote Guardian" — search for "remote" in Settings to find it.
  \item \textbf{Universal Solution (Works for All Brands)}: Download the "Sunflower Remote Control" app (free version is sufficient). Install it on both phones. Register the same account or add as friends. Have your parents read you the verification code displayed on their phone, and you can remotely view and operate their screen.
  \item \textbf{: If you need to extract text from scanned documents for editing, you need OCR. WeChat's built-in "Extract Text" feature is the simplest — long-press the image to use it, and the recognition accuracy is decent. For bulk processing, you can use the Baimiao app or CamScanner. For Chinese recognition, Baidu's OCR engine works quite well; WPS's image-to-text feature is also very accurate.}: FaceTime supports "Screen Sharing" during calls. While on a FaceTime call, tap the screen share button, and the other person can see your screen. But iOS security prevents remote operation — it's "instructive teaching" only.
\end{itemize}

\booksection{V. WeChat “Senior-Friendly Makeover” — Just Tell Your Parents These Things}

WeChat is the app elderly people use most, but the interface is too dense and cluttered. Here are some quick, effective adjustments:

\begin{itemize}
  \item \textbf{Enlarge WeChat's Own Font}: WeChat > Me > Settings > General > Font Size, drag to maximum. Note: this only affects the font within WeChat, not the system font.
  \item \textbf{Enable "Care Mode"}: WeChat > Me > Settings > Care Mode > Enable. After enabling, fonts and buttons become larger, color contrast is enhanced, and you can one-tap "Listen to Text Messages" (tap a text message and it reads aloud).
  \item \textbf{Help Them Label Important Contacts}: Change children's and relatives' names in contacts to easily recognizable ones — like "Eldest Son - Xiaoming" or "Second Daughter - Xiaohong," so they immediately know who's who when opening WeChat.
  \item \textbf{Teach Them to Use "Voice Input" Instead of Typing}: On the keyboard next to the chat box, hold the spacebar (Android) or tap the microphone icon (iPhone), and speak directly — it converts to text and sends. Recognition accuracy is already very high — for elderly people who can't type in pinyin, it's a "game-changing" convenience.
  \item \textbf{Disable Unused Features in "Discover"}: WeChat > Me > Settings > General > Discover Page Management, turn off "Moments," "Channels," "Live," "Shopping," "Games" — anything not used — to reduce clutter.
\end{itemize}

\booksection{VI. What If the Phone Is Lost? Make an “Emergency Plan” for Elderly Relatives}

Preparation prevents problems. While the phone isn't lost yet, do these things first:

\begin{enumerate}
  \item \textbf{Confirm "Find My Phone" Is Enabled}: iPhone — Settings > [Your Name] > Find My > Find My iPhone; Android — search for "Find Device" in Settings. Make sure this toggle is on.
  \item \textbf{Write Important Passwords and Accounts on a Piece of Paper}: Lock one copy of Apple ID / Huawei account / Xiaomi account passwords in a home drawer, and give another copy to a trusted child. If the phone is lost, you need these accounts to remotely lock and locate it — forgetting the password makes it very difficult.
  \item \textbf{Teach Elderly Relatives One Sentence}:\textbf{"If the phone is lost, the first thing to do is call family."}Family can immediately log into iCloud or Android's Find My Device website to remotely mark it as "Lost Mode" — the person who finds it can't unlock or use it.
\end{enumerate}

\booksection{VII. Wireless Earphones and Hearing Aid Compatibility — An Overlooked Feature}

Many elderly relatives have hearing issues but don't know that phones can work with hearing aids. iPhone and some Android phones support MFi (Made for iPhone) hearing aids, which can directly transmit phone call audio to the hearing aid — much clearer than holding the phone to the ear.

\begin{itemize}
  \item \textbf{Connection Method}: iPhone Settings > Accessibility > Hearing Devices, turn on Bluetooth, and follow the hearing aid's pairing instructions.
  \item \textbf{If Not Using a Hearing Aid}: In "Accessibility > Audio/Visual," enable "Phone Noise Cancellation" and "Mono Audio" to improve call clarity.
  \item \textbf{Android}: In "Accessibility > Hearing Enhancement," there are similar settings — locations vary by brand. Search for "hearing aid" directly.
\end{itemize}

\booksection{VIII. Health-Related Features — A Phone Is More Than a Communication Tool}

\begin{itemize}
  \item \textbf{SOS Emergency}: Press the power button five times consecutively (on most phones), and the phone automatically calls emergency services and sends location to preset emergency contacts. Set this up for elderly relatives in advance — iPhone: "Settings > Emergency SOS"; Android: "Settings > Safety > Emergency Assistance." Practice once before teaching them, so they don't fumble when stressed.
  \item \textbf{Fall Detection}: Apple Watch has this feature — if it detects a fall followed by no movement, it automatically calls emergency services and notifies family. If budget allows, getting one for elderly relatives can save lives at critical moments.
  \item \textbf{Medication Reminders}: iPhone's "Health" app has a built-in medication reminder — set the medication name, time, and dosage, and it pops up a reminder when it's time. Android can use "alarm clocks" or download "medication reminder" apps.
\end{itemize}

\begin{tcbsummary}{Chapter Summary}
\begin{itemize}
  \item Enable Easy Mode or manually enlarge fonts/icons first — once the environment is set up, most problems disappear
  \item Three iron rules for helping elderly relatives with phones: don't rush, use analogies, write it down
  \item Remind about six common scams repeatedly — three defense lines can block 90\% of fraud
  \item Remote assistance lets you fix your parents' phone from anywhere
  \item WeChat's "Care Mode" + "Voice Input" are the two most practical features for elderly users
  \item SOS Emergency and Fall Detection — life-saving configurations worth spending five minutes to set up
\end{itemize}

\end{tcbsummary}
